\subsection{Channel-centre vertical energy bands}
\begin{figure}[H]\centering
\begin{tikzpicture}
\begin{axis}[width=.92\linewidth,height=6.8cm,
xlabel={Depth from gate-side silicon interface (nm)},
ylabel={$E_C$ (eV)},xmin=0,xmax=30,
ymin=.375,ymax=.47,grid=both,legend pos=north west]
\addplot+[no marks,thick]
table[col sep=comma,x=Depth_nm,y=Ec_eV]
{data/part5_vertical_tsi30.csv};
\addlegendentry{$T_{\mathrm{SI}}=30$ nm}
\addplot+[no marks,thick,dashed]
table[col sep=comma,x=Depth_nm,y=Ec_eV]
{data/part5_vertical_tsi10.csv};
\addlegendentry{$T_{\mathrm{SI}}=10$ nm}
\end{axis}
\end{tikzpicture}
\caption{Channel-centre vertical conduction-band energy
at $V_{GS}=0$ V and $V_{DS}=1.2$ V.
The 10-nm trace ends at the 10-nm Si/BOX interface;
the 30-nm trace ends at 30 nm. Only silicon is shown.}
\label{fig:part5_vert_ec}
\end{figure}

\begin{figure}[H]\centering
\begin{tikzpicture}
\begin{axis}[width=.92\linewidth,height=6.8cm,
xlabel={Depth from gate-side silicon interface (nm)},
ylabel={$E_V$ (eV)},xmin=0,xmax=30,
ymin=-.735,ymax=-.635,grid=both,legend pos=north west]
\addplot+[no marks,thick]
table[col sep=comma,x=Depth_nm,y=Ev_eV]
{data/part5_vertical_tsi30.csv};
\addlegendentry{$T_{\mathrm{SI}}=30$ nm}
\addplot+[no marks,thick,dashed]
table[col sep=comma,x=Depth_nm,y=Ev_eV]
{data/part5_vertical_tsi10.csv};
\addlegendentry{$T_{\mathrm{SI}}=10$ nm}
\end{axis}
\end{tikzpicture}
\caption{Channel-centre vertical valence-band energy,
using the same off-gate/high-drain operating points
as Fig.~\ref{fig:part5_vert_ec}.}
\label{fig:part5_vert_ev}
\end{figure}

\paragraph{Vertical bending and gate control.}
The 30-nm device exhibits an increase in
$E_C$ from 0.39347 eV near the gate to
0.45855 eV near the buried oxide:
a total of 65.1 meV through the top
silicon. The corresponding increase for
10-nm silicon is 0.38269 to 0.39275 eV,
or just 10.1 meV.
The valence bands show similar changes
of 64.9 meV and 10.4 meV, respectively.
Thus the thinner silicon film displays a more
nearly uniform transverse electrostatic band
profile, indicating reduced variation of
potential across the body. This is consistent
with improved geometric gate control and
the 10-nm device's reduced OFF current
observed in Part~2. Both energy bands
shift together within each film, with a local
silicon bandgap of about 1.105 eV.

\paragraph{Fully depleted versus partially depleted.}
The original vertical SDevice records, before
removing auxiliary columns for the report plots,
contain mobile-hole densities far below the
local acceptor densities throughout both films.
The maximum ratios of hole density to acceptor
concentration were $7.5\times10^{-8}$
(30 nm) and $5.8\times10^{-8}$ (10 nm).
The recorded net space charge was also
approximately the negative acceptor density
across the entire channel-centre silicon,
within about 0.01\% in these samples.
There is therefore no observed neutral,
hole-rich core along either cutline:
\emph{both devices are consistent with full
body depletion at $V_{GS}=0$ V,
$V_{DS}=1.2$ V}. The larger 30-nm band
bending is not by itself evidence of
partial depletion. The given data do not
establish full depletion at all terminal
biases, but do directly support depleted
bodies at the specified Part~5 condition.

\paragraph{Evidence limitations.}
The archived TDRs were used without rerunning
the physics. The four figures are based on
Sentaurus numerical samples, not idealized
band sketches. The longitudinal plot
subsamples reproduce the original raw
traces to within 2 meV under linear
interpolation; vertical silicon samples
were not subsampled. These band-edge
snapshots cannot independently produce a
numerical DIBL value without low-drain-bias
band profiles; any stronger claim would
exceed the evidence obtained here.
